%% ma: L10-B13-0053
%% lop: 10
%% chuong: 5
%% bai: 13
%% dang: 10.13.7
%% loai: TLN
%% muc: TH
%% dap_an: "4"
%% nguon: "chuyen-de-cac-so-dac-trung-cua-mau-so-lieu-khong-ghep-nhom-toan-10-chuong-trinh-moi.pdf (GV. Phan Nhật Linh) (P02), Bài 02 Dạng 1 Phần III Câu 3"
%% kien_thuc: "Lớp 10 Bài 13 (đọc và phân tích số liệu), tỉ lệ phần trăm"
%% trang_thai: chinh-thuc
%% ghi_chu: "Tự duyệt theo yêu cầu giáo viên 22/10/2026."
\begin{ex}
Trong $7$ tháng đầu năm, số sản phẩm sản xuất mỗi tháng của công ty X đều tăng khoảng $5\%$ so với tháng trước đó. Biết rằng trong bảng dưới đây, số sản phẩm của đúng một tháng bị nhập sai. Tháng đó là tháng thứ mấy?
\begin{center}
\begin{tabular}{|c|c|c|c|c|c|c|c|}\hline
Tháng & $1$ & $2$ & $3$ & $4$ & $5$ & $6$ & $7$\\\hline
Số sản phẩm & $500$ & $525$ & $551$ & $569$ & $606$ & $636$ & $668$\\\hline
\end{tabular}
\end{center}
\shortans{4}
\loigiai{Tỉ lệ tăng so với tháng trước: $5\%$; $4{,}95\%$; $3{,}27\%$; $6{,}50\%$; $4{,}95\%$; $5{,}03\%$.

Chỉ tháng 4 và tháng 5 lệch khỏi $5\%$, mà $551\cdot1{,}05\approx579$ và $579\cdot1{,}05\approx608\approx606$, nên số liệu tháng $4$ bị nhập sai.}
\end{ex}
